\begin{center}
\begin{minipage}{0.98\linewidth}
\begin{tcolorbox}[colback=white,colframe=CStatement!25,boxrule=0.4pt,
  arc=1mm,left=3mm,right=3mm,top=3mm,bottom=3mm]
\centering
\textbf{根号来自斜率方向三角形}\par\medskip
\begin{tikzpicture}[font=\small,>=Stealth]

\coordinate (U) at (0,0);\coordinate (V) at (4,0);\coordinate (W) at (4,3);
\draw[CStatement,thick] (U)--(V)--(W)--cycle;
\draw (3.75,0)--(3.75,.25)--(4,.25);
\pic[draw,angle radius=10mm,angle eccentricity=1.4,"$\theta$"] {angle=V--U--W};
\node[below] at (2,0) {$1$};
\node[right] at (4,1.5) {$|k|$};
\node[above left,rotate=36.87] at (2,1.5) {$\sqrt{1+k^2}$};
\node[align=left,anchor=west] at (5.25,1.5) {$\displaystyle\cos\theta=\frac{1}{\sqrt{1+k^2}}$};

\end{tikzpicture}\par\medskip
图按 $|k|=3/4$ 等比例绘制。把水平边长记为 $1$，竖直边长就是 $|k|$；由勾股定理，斜边长为 $\sqrt{1+k^2}$。
\end{tcolorbox}
\end{minipage}
\end{center}
